\begin{textsample}{POS Dim 3 – human – Score 10.00 – mg/mg\_ef\_linguagens\_ed\_fis\_3\_p4.txt}  \label{ex:f3_pos_001}
Com base nas reflexões anteriores sobre Educação e Educação Física, nos eixos norteadores, tanto das \textbf{diretrizes} \textbf{curriculares} \textbf{nacionais} para o Ensino Fundamental e Médio como das \textbf{diretrizes} \textbf{curriculares} \textbf{propostas} para a formação de professores da educação básica, iremos discutir alguns princípios que julgamos fundamentais para orientar as ações \textbf{educativas} e os processos de tomada de decisões dos educadores, em especial no que se refere à Educação Física nas Séries Finais do Ensino Fundamental.19 19 Os CBCs foram construídos para nortear a organização \textbf{curriculares} dos anos finais do ensino fundamental. \textbf{Cabe} destacar que neste momento de construção de um \textbf{documento} que abrange toda esta \textbf{etapa} ( anos iniciais e anos finais ) os pressupostos aqui apresentados são pertinentes também para todos esses anos e as \textbf{especificidades} relativas aos sujeitos atendidos ( especialmente em se tratando dos \textbf{primeiros} anos do ensino fundamental – tempo de transição da infância ). É importante refletir sobre estratégias e metodologias que acolham e respeitem o que é singular em cada tempo de desenvolvimento humano.

% matched lemmas: caber, curricular, diretriz, documento, educativo, especificidade, etapa, nacional, primeiro, proposta
\end{textsample}
